\section{Taller resuelto: Leer un mapa adimensional de hélice}

\noindent\textbf{Nivel:} Intermedio.

\apartado{1. Enunciado}{Una hélice de diámetro 2.1 m gira a 32 rev/s con avance 58 m/s. En ese punto se han dado CT=0.083 y CP=0.090, con densidad 1.08 kg/m³. Calcule J, empuje, potencia al eje y rendimiento.}

\apartado{2. Datos}{Los coeficientes pertenecen al mismo punto de mapa y se usan las definiciones de este capítulo.}

\apartado{3. Incógnitas}{Las magnitudes y comprobaciones solicitadas en el caso.}

\apartado{4. Sistema de referencia}{Empuje paralelo al avance aéreo; coeficientes de hélice adimensionales.}

\apartado{5. Hipótesis físicas}{Régimen y paso fijos en el punto; los valores de mapa son datos docentes, no mediciones.}

\apartado{6. Modelo matemático}{Leer un mapa adimensional de hélice}

\apartado{7. Ecuaciones simbólicas}{Se parte de los balances y definiciones siguientes.}

\begin{equation}J=V/(n_rd)\end{equation}

\begin{equation}T=C_T\rho n_r^2d^4\end{equation}

\begin{equation}P=C_P\rho n_r^3d^5\end{equation}

\begin{equation}\eta_p=JC_T/C_P\end{equation}

\apartado{8. Desarrollo algebraico}{Las operaciones se realizan en este orden, conservando las condiciones del modelo.}

\begin{enumerate}

\item Se usa régimen en revoluciones por segundo, sin multiplicar por 2pi en estas definiciones.

\item Se calculan empuje y potencia con las potencias correctas de diámetro.

\item Se obtiene eficiencia por coeficientes y se comprueba con TV/P.

\end{enumerate}

\apartado{9. Sustitución numérica}{Se conservan los resultados intermedios sin redondear y se muestran todas las operaciones escalares disponibles.}

\begin{enumerate}

\item Avance J.

\[\fitmath{y_{1}=\frac{58}{32\cdot 2.1}\approx 0.8630952\,}\]

\item Empuje.

\[\fitmath{y_{2}=0.083\cdot 1.08\cdot {32}^{2}\cdot {2.1}^{4}\approx 1785.168\,\mathrm{N}}\]

\item Potencia al eje.

\[\fitmath{y_{3}=0.09\cdot 1.08\cdot {32}^{3}\cdot {2.1}^{5}\approx 130080.6\,\mathrm{W}}\]

\item Rendimiento.

\[\fitmath{y_{4}=\frac{\frac{58}{32\cdot 2.1}\cdot 0.083}{0.09}\approx 0.7959656\,}\]

\end{enumerate}


\apartado{10. Resultados con unidades}{Los resultados pertenecen al modelo y al intervalo de datos declarados.}

\begin{center}\begin{tabular}{lr}\toprule Magnitud&Resultado\\\midrule

Avance J & $0.8630952\,$\\

Empuje & $1785.168\,\mathrm N$\\

Potencia al eje & $130080.6\,\mathrm W$\\

Rendimiento & $0.7959656\,$\\

\bottomrule\end{tabular}\end{center}

\apartado{11. Interpretación}{Un mapa convierte la hipótesis de rendimiento constante en una función de condición; no se extrapola este punto al reposo.}

\apartado{12. Comprobación independiente y alcance}{Calcular TV/P con fuerzas y potencia dimensionales y compararlo con J CT/CP.}

\apartado{13. Limitaciones}{No se han comprobado Mach de punta ni límites de par del motor; un rendimiento mayor que uno indicaría datos o convenciones incompatibles.}

\par\smallskip\noindent{\footnotesize\textit{Datos y solución originales. Las sustituciones registradas se reproducen con SymPy y Studium; eso no establece validación experimental ni revisión experta.}}
